% Copyright (c) 2026 Eugene Joseph M. Ragasa.
% SPDX-License-Identifier: Apache-2.0
% Author: Eugene Joseph M. Ragasa
% Formal proof completed: 2026-10-20
\documentclass[11pt]{article}
\usepackage[T1]{fontenc}
\usepackage{amsmath,amssymb,booktabs,tabularx}
\usepackage{listings}
\usepackage[margin=1in]{geometry}
\usepackage{xurl}
\usepackage[hidelinks]{hyperref}
\setlength{\emergencystretch}{3em}
\lstset{basicstyle=\ttfamily\small,breaklines=true,columns=fullflexible}
\title{Proof Dossier: Algebraic Specification of the Lower-Band Projector}
\author{Eugene Joseph M. Ragasa}
\date{Formal proof completed: 2026-10-20}
\begin{document}\sloppy\maketitle
\section{Identity and status}
\begin{tabularx}{\textwidth}{@{}lX@{}}\toprule
Statement ID & \path{BG-PROJECTION-THM-005} \\
Status & proved; no admissions \\
Lean declaration & \path{ProjectKoios.BandGauge.Projection.twoBandLowerBandProjector_specification} \\
Owning module & \path{ProjectKoios/BandGauge/Projection/SpectralProjectionRange.lean} \\
Axioms & \texttt{propext}, \texttt{Classical.choice}, \texttt{Quot.sound} \\
\bottomrule\end{tabularx}

\section{Exact mathematical statement}
For the source parent $H(k)$, positive dispersion $d(k)$, and $\Delta t\ne0$, define
\[
 P_-(k)=\frac12\left(I-\frac{H(k)}{d(k)}\right).
\]
Then, for every real $k$,
\[
 P_-(k)^2=P_-(k),\qquad H(k)P_-(k)=-d(k)P_-(k).
\]

\section{Source provenance and relationship}
This theorem \emph{derives} the algebraic projector properties implicit in the lower-band
notation of \path{manuscript.tex:324--334} at owner commit
\path{99a46cb59d0c6f27ada7d709f194996ef9619d9f}, manuscript SHA-256
\texttt{d8d5951d8d7ed75355b5d94f37fcbddc70c1466897a5770dbe9f877e60016daf}.
The cited source bytes are clean at that immutable commit.

\section{Assumptions, conventions, and delta}
The source assumption $\Delta t\ne0$ is preserved, implying $d(k)>0$ and making $d^{-1}$
well defined. Matrices use the displayed source basis. The theorem records idempotency and
the lower eigenrelation explicitly; these are consequences of the source parent-square
identity, not an additional physical claim.

\section{Statement-specific prior art}
\begin{tabularx}{\textwidth}{@{}lXXX@{}}\toprule
Relation & Identity & Locator & Role \\
\midrule
derives & owner source above & manuscript lines 324--334 & justifies the symbol $P_-$ \\
depends-on & \path{BG-PROJECTION-THM-001} & theorem manifest & supplies $H^2=d^2I$ after circle evaluation \\
background-only & canonical external identity unresolved & unverified & standard two-eigenvalue projector formula; no novelty claim \\
\bottomrule\end{tabularx}

\section{Motivation and contribution boundary}
The theorem prevents the phrase ``lower-band projector'' from resting only on a name. It
checks idempotency and places the formula's range in the $-d$ eigenspace. It does not
separately prove self-adjointness, establish a general spectral theorem, or bind a material
model.

\section{Proof schematic and argument}
From $H^2=d^2I$ and $d\ne0$, expand
\[
 P_-^2=\tfrac14(I-2H/d+H^2/d^2)=\tfrac12(I-H/d)=P_-.
\]
Similarly,
\[
 HP_-=\tfrac12(H-H^2/d)=\tfrac12(H-dI)=-dP_-.
\]
The positivity of $d$ follows from the nonzero gap term in its defining square root.

\section{Lean realization}
Private matrix-algebra lemmas prove idempotency and the eigenrelation from a generic
rank-two operator-square hypothesis. The source theorem instantiates them with
\path{twoBandParentOnUnitCircle_sq}. Scalar denominators are discharged with exact
field normalization. Toolchain: Lean 4.35.0, Mathlib
\texttt{56a005855ea17e0b0958f73ec801ae140b0a1679}.

\section{Trusted base and axiom report}
\begin{lstlisting}
#print axioms ProjectKoios.BandGauge.Projection.twoBandLowerBandProjector_specification
-- [propext, Classical.choice, Quot.sound]
\end{lstlisting}

\section{AI, automation, and comprehension record}
The implementation and prose were AI-assisted. Lean kernel checking validates the exact
conjunction. Field simplification and ring normalization produced checked proof terms.
The maintainer reviewed both displayed expansions. No independent expert review has
occurred.

\section{Scientific and evidence boundaries}
The theorem is pointwise exact matrix algebra. It does not prove locality, decay, a
numerical eigensystem, physical validity, or a universal statement about projectors.

\section{Unresolved questions}
A separate Hermiticity/self-adjointness declaration could support a more general abstract
orthogonal-projector API, but is not needed for the exact source sandwich identity.
Canonical external prior-art identity remains unresolved.

\section{Review and revision history}
\begin{tabularx}{\textwidth}{@{}lllX@{}}\toprule
Date & Revision & Process & Outcome \\
\midrule
2026-10-20 & initial & AI-assisted implementation; Lean kernel check & projector idempotency and lower-eigenrelation proved; no independent expert review \\
\bottomrule\end{tabularx}
\end{document}
